\documentclass[10pt,a4paper]{amsart}
\usepackage[margin=19mm]{geometry}
\usepackage{amsmath,amssymb,mathtools}
\usepackage[scr=rsfs,cal=euler]{mathalfa}
\usepackage{xfrac}
\usepackage[unicode,hidelinks,bookmarks=false]{hyperref}
\newcommand{\ie}{i.e.\ }
\newcommand{\cf}{cf.\ }
\newcommand{\resp}{respectively\ }
\newcommand{\Subsection}{\subsection}
\providecommand{\nobreakdash}{\nobreak\hbox{-}}
\setlength{\emergencystretch}{2em}
\multlinegap0pt
\usepackage{mathtools}
\newtheorem{theorem}{Theorem}[section]
\newtheorem{proposition}[theorem]{Proposition}
\newtheorem{lemma}[theorem]{Lemma}
\newtheorem{corollary}[theorem]{Corollary}
\newtheorem{definition}[theorem]{Definition}
\newtheorem{example}[theorem]{Example}
\newtheorem{remark}[theorem]{Remark}
\newcommand{\sa}{\mathrm{sa}}
\newcommand{\QM}{\mathrm{QM}}
\newcommand{\LM}{\mathrm{LM}}
\newcommand{\RM}{\mathrm{RM}}
\newcommand{\M}{\mathrm{M}}
\begin{document}
\title[Order isomorphisms on positive cones of non-unital C^*-algeb]{Order isomorphisms on positive cones of non-unital $C^*$-algebras}
\author{Osamu Hatori; Shiho Oi}
\date{}
\maketitle
\noindent\emph{Source and licence.} Order isomorphisms on positive cones of non-unital $C^*$-algebras. Version arXiv:2609.21339v1 (CC BY 4.0). This material is distributed under the Creative Commons Attribution 4.0 International licence, http://creativecommons.org/licenses/by/4.0/, as recorded in the arXiv OAI licence metadata for this exact version and shown on the abstract page https://arxiv.org/abs/2609.21339v1. Full rights notice: licenses/arxiv-2609.21339-CC-BY-4.0.txt.\par\smallskip
\noindent\textit{Editorial scope.} Counted unit: the original \texttt{theorem} environment \texttt{thm:extension} at source lines 170--177 together with its complete proof at lines 179--271; the delivered document keeps source lines 97--272 so that every referenced label and every citation resolves inside it. Native editable source is the paper's own concatenated TeX; the AMS wrapper is editorial formatting only. No publisher class, upstream script or unrelated artwork is redistributed.
\section{Positively homogeneous order isomorphisms}\label{section2}
We study positively homogeneous order isomorphisms on the positive cones of $C^{*}$-algebras.  The key input is Sch\"affer's order-unit linearization theorem \cite[Theorem~B]{SchafferGauge}, in the
form recalled in \cite[Theorem~5.2]{LGI}.
Fix $u\in A_+\setminus\{0\}$ and put
\[
 X_u=\{x\in A_{\sa}: -\lambda u\le x\le \lambda u
       \text{ for some }\lambda>0\},
 \qquad
 C_u=X_u\cap A_+.
\]
Then $X_u$ is a real partially ordered vector space with positive cone
$C_u$.  By the definition of $X_u$, the element $u$ is an order unit.
Moreover, $(X_u,C_u,u)$ is Archimedean.  Indeed, suppose that $x\in X_u$
and $y\in C_u$ satisfy $nx\le y$ for every $n\ge1$.  Since
$y\le\|y\|1_{A^{**}}$, we have
\[
        nx\le\|y\|1_{A^{**}} \qquad(n\ge1).
\]
If $x\not\le0$, then $\sup\sigma(x)>0$, and hence
$n\sup\sigma(x)\le\|y\|$ for every $n$, a contradiction.  Thus $x\le0$.
We write
\[
 \|x\|_u=\inf\{\lambda>0:-\lambda u\le x\le\lambda u\}
\]
for the corresponding order-unit norm.

\begin{lemma}\label{lem:interior}
The interior of $C_u$ with respect to the order-unit norm is
\[
 C_u^\circ
 =\{x\in A_+: \alpha u\le x\le\beta u
       \text{ for some }\alpha,\beta>0\}.
\]
\end{lemma}

\begin{proof}
Suppose first that $\alpha u\le x\le\beta u$ for some
$\alpha,\beta>0$.  If $\|y-x\|_u<\alpha/2$, choose
$\delta<\alpha/2$ such that
\[
        -\delta u\le y-x\le\delta u.
\]
Then $y\ge(\alpha-\delta)u\ge(\alpha/2)u\ge0$.  Thus
$x\in C_u^\circ$.

Conversely, if $x\in C_u^\circ$, then there is $\varepsilon>0$ such
that $x+y\in C_u$ whenever $\|y\|_u<\varepsilon$.  Taking
$y=-(\varepsilon/2)u$ gives
\[
 \frac\varepsilon2u\le x.
\]
Since $x\in X_u$, there is $\beta>0$ with $x\le\beta u$.
\end{proof}

\begin{lemma}\label{lem:span-interior}
We have $X_u=\operatorname{span}_{\mathbb R} C_u^\circ$. 
\end{lemma}

\begin{proof}
Let $x\in X_u$.  Choose $\lambda>0$ with
$-\lambda u\le x\le\lambda u$.  Then
\[
 u\le x+(\lambda+1)u\le(2\lambda+1)u,
\]
so $x+(\lambda+1)u\in C_u^\circ$ by
Lemma~\ref{lem:interior}.  Also $(\lambda+1)u\in C_u^\circ$, whence
\[
 x=\bigl(x+(\lambda+1)u\bigr)-(\lambda+1)u
 \in\operatorname{span}_{\mathbb R}C_u^\circ.
\]
The reverse inclusion is immediate.
\end{proof}


\begin{theorem}\label{thm:extension}
Let $A$ and $B$ be $C^*$-algebras and let $ T: A_+\to B_+$ be a positively homogeneous order isomorphism. 
Then $T$ is additive on $A_+$ and extends uniquely to a real-linear
order isomorphism
\[
 \widetilde T:A_{\sa}\longrightarrow B_{\sa}.
\]
\end{theorem}

\begin{proof}
Since $0$ is the least element of each positive cone and $T$ is an
order isomorphism, $T(0)=0$.  In particular, $u\ne0$ implies
$T(u)\ne0$.

Fix $u\in A_+\setminus\{0\}$.  By Lemma~\ref{lem:interior}, since $T$ and $T^{-1}$ are 
order-preserving and positively homogeneous, 
\[
 T(C_u^\circ)=C_{T(u)}^\circ.
\]
Hence
\[
 \Phi_u:=T|_{C_u^\circ}:C_u^\circ\longrightarrow C_{T(u)}^\circ
\]
is a positively homogeneous order isomorphism.  By Sch\"affer's
order-unit linearization theorem, in the form recalled in
\cite[Theorem~5.2]{LGI}, $\Phi_u$ is linear in the sense that it is the
restriction of a real-linear map on the span of $C_u^\circ$.  Since
$C_u^\circ$ and $C_{T(u)}^\circ$ span $X_u$ and $X_{T(u)}$, respectively,
by Lemma~\ref{lem:span-interior}, $\Phi_u$ extends uniquely to a
real-linear map
\[
 \Phi_u:X_u\longrightarrow X_{T(u)}.
\]
We claim that $\Phi_u(x)=T(x)$ for every $x\in C_u$.  
Fix $x\in C_u$.  We first prove that $\Phi_u(x)\in B_+$. As $x\in C_u$, $x+\varepsilon u\in C_u^\circ$ for every $\varepsilon>0$, and hence
\[
        \Phi_u(x)+\varepsilon\Phi_u(u)
        =\Phi_u(x+\varepsilon u)\in C_{T(u)}\subset B_+.
\]
Letting $\varepsilon\downarrow0$ and using the norm-closedness of $B_{+}$, we obtain $\Phi_u(x)\in B_+$. 
Next we prove $T(x)\le \Phi_u(x)$. 
%Let $x\in C_u$.
Choose $\lambda>0$ with $0\le x\le\lambda u$.  For every
$\varepsilon>0$,
\[
 \varepsilon u\le x+\varepsilon u
 \le(\lambda+\varepsilon)u,
\]
so $x+\varepsilon u\in C_u^\circ$.  Therefore
\[
 T(x)\le T(x+\varepsilon u)
 =\Phi_u(x)+\varepsilon\Phi_u(u).
\]
Letting $\varepsilon\downarrow0$ and using the norm-closedness of $B_{+}$, we obtain 
$T(x)\le\Phi_u(x)$.


Since $\Phi_u(x)\in B_+$, we may apply $T^{-1}$.
From
\[
    0\le\Phi_u(x)
    \leq\Phi_u(x)+\varepsilon\Phi_u(u)
    =T(x+\varepsilon u)
\]
and the order preservation of $T^{-1}$, we obtain
\[
    T^{-1}(\Phi_u(x))\leq x+\varepsilon u
    \qquad(\varepsilon>0).
\]
Again letting $\varepsilon\downarrow0$ gives
$T^{-1}(\Phi_u(x))\le x$, and hence $\Phi_u(x)\le T(x)$.  Thus $T(x)=\Phi_u(x)$ for all $x\in C_u$.

Now let $a,b\in A_+$, not both zero, and put $u=a+b$.  Since
$a,b,a+b\in C_u$,
\[
 T(a+b)=\Phi_u(a+b)=\Phi_u(a)+\Phi_u(b)=T(a)+T(b).
\]
The case $a=b=0$ is trivial, so $T$ is additive on $A_+$.

For $x\in A_{\sa}$ write $x=a_1-a_2$ with $a_1,a_2\in A_+$ and set
\[
 \widetilde T(x)=T(a_1)-T(a_2).
\]
If $a_1-a_2=b_1-b_2$, then $a_1+b_2=b_1+a_2$, and additivity gives
\[
 T(a_1)+T(b_2)=T(b_1)+T(a_2),
\]
so the definition is independent of the chosen decomposition.
Additivity and positive homogeneity of $T$ imply real linearity of
$\widetilde T$.  The same construction applied to $T^{-1}$ shows that
$\widetilde T$ is bijective.

Finally, if $x\le y$, then $y-x\in A_+$ and therefore
\[
 \widetilde T(y)-\widetilde T(x)
 =\widetilde T(y-x)=T(y-x)\ge0.
\]
The same argument for $\widetilde T^{-1}$ shows that the order is
preserved in both directions.  Uniqueness follows because
$A_{\sa}=A_+-A_+$, so a real-linear extension is determined by its
values on $A_+$.
\end{proof}

\begin{thebibliography}{99}
\bibitem[LGI]{LGI} B.~Lemmens, O.~van~Gaans and H.~van~Imhoff, \textit{On the linearity of order-isomorphisms}, Canad. J. Math. \textbf{73} (2021), no.~2, 399--416.
\bibitem[SchafferGauge]{SchafferGauge} J.~J. Sch\"affer, \textit{Orders, gauge, and distance in faceless linear cones. II. Gauge-preserving bijections are cone-isomorphisms}, Arch. Rational Mech. Anal. \textbf{67} (1978), 305--313.
\end{thebibliography}
\end{document}
